\chapter{Results and Analysis}
\setcounter{equation}{0}
\section{Feature Stability and Dataset Differences}
The article compares feature importance across the custom, GramBeddings, and PhiUSIIL datasets. URL length was identified as a relatively robust feature, with reported mean importance of 0.0319 and coefficient of variation of 0.7843. HTTPS showed much greater variation in importance, from 0.125 in GramBeddings to 0.513 in PhiUSIIL. This is consistent with the observed dataset-dependent effect of HTTPS on accuracy.\cite{bib1}

\section{Why Dataset Shift Matters}
Datasets can differ in URL length distributions, HTTPS adoption, phishing prevalence, collection period, and source. A classifier may inadvertently learn patterns that distinguish the labels in one dataset but do not persist in another. Cross-dataset tests reveal this risk and give a more realistic view of the model's limits than a single random split.

\section{Evaluation Metrics}
Accuracy measures the proportion of all samples classified correctly. Recall measures the proportion of phishing URLs detected. Specificity measures the proportion of legitimate URLs correctly identified. Precision measures how many URLs flagged as phishing are actually phishing, and the F1 score balances precision and recall. No single metric fully describes operational performance, so the paper reports several metrics and evaluates multiple datasets.\cite{bib1}

\section{Interpretation}
The evidence supports three conclusions. First, the selected features can support strong detection on the custom dataset. Second, performance declines on independent datasets, meaning the detector requires continued evaluation after deployment. Third, feature importance is contextual: a signal such as HTTPS may be useful in one dataset but nearly uninformative in another. The results do not imply that all future URL populations will achieve the same scores.
